\section{Notation and Conventions}

We use energy and distance interchangeably: lower values indicate better
matches. For an input and $K$ components, $d_j$ denotes the energy of component
$j$, and
\[
    r = \operatorname{softmax}(-d)
\]
denotes the responsibility vector. Thus $r_j \ge 0$ and $\sum_j r_j = 1$.

The log-sum-exp objective is
\[
    L_{\mathrm{LSE}}(d) = -\log \sum_j \exp(-d_j),
\]
so that
\[
    \frac{\partial L_{\mathrm{LSE}}}{\partial d_j} = r_j.
\]
This sign convention is fixed throughout the paper.

For supervised classification, $h$ denotes logits and
$p=\operatorname{softmax}(h)$ denotes the output class posterior. We write
$W_1$ for the input-to-component map and $W_2$ for the learned dense
component-to-class map. The latter is the composition corridor whose effect on
conditioning is the focus of the paper.

The NegLogSoftmin calibration is
\[
    \operatorname{NLS}(d)_j
    = d_j + \log \sum_k \exp(-d_k),
\]
which preserves relative distances and satisfies
\[
    \exp(-\operatorname{NLS}(d)_j) = r_j.
\]

Volume control consists of an anti-collapse variance penalty $L_{\mathrm{var}}$
and an anti-redundancy decorrelation or total-correlation penalty
$L_{\mathrm{tc}}$. Unless otherwise stated, $\lambda$ denotes the total
auxiliary weight applied to the local EM objective.
